%=====================================================
%
% Variables (same scheme as the sample document)
%
%=====================================================

\def\doctype{article}

\def\emptypages{false}
\def\showparts{false}

\def\ttitlepage{titlepage.tex}
\def\ttitlepagetitletop{Planner and Worker}
\def\ttitlepagesubtitletop{When a Cloud Planner over a Local Coder Pays Off}
\def\ttitlepagetitlemiddle{Lab Report}
\def\ttitlepagesubtitlemiddle{Measurements of 15--16 November 2026}

\def\torganization{Independent Research}
\def\tsubmissiondate{16 November 2026}
\def\tversion{0.1 (draft)}
\def\tauthor{Matthias Nott}
\def\tcontact{~}
\def\tcontactid{~}

\def\tfooterleftlabelup{Title}
\def\tfooterleftlabeldown{Author}
\def\tfooterleftup{Local Coding Agents, Report 4}
\def\tfooterleftdown{Matthias Nott}
\def\tfootermiddleup{Planner and Worker}
\def\tfootermiddledown{Lab Report}
\def\tfooterrightlabelup{Date}
\def\tfooterrightlabeldown{Version}
\def\tfooterrightup{November 2026}
\def\tfooterrightdown{0.1}

\def\showtitlepage{true}
\def\showlistofs{true}
\def\showtoc{true}
\def\showlot{true}
\def\showlof{false}
\def\showloe{false}
\def\showlod{false}

\def\showreferences{true}
\def\showappendix{false}
\def\showendnotes{false}
\def\showanyglossary{false}
\def\showacronyms{false}
\def\showglossary{false}
\def\showstatic{false}
\def\showindex{false}

\def\showpageheaders{true}
\def\showpagefooters{true}
\def\showwordcount{false}
\def\showwordcountback{false}

\input{templates/cfg/zz_layout}

\setcounter{tocdepth}{2}

\usepackage[
  style=apa
, citestyle=authoryear
, backend=biber
, maxnames=3
, minnames=1
, sorting=nyt
, sortcites=true
, block=space
, safeinputenc
, natbib=true
, backref=true
, uniquename=init]{biblatex}
\DeclareLanguageMapping{american}{american-apa}
\addbibresource{Bibliography.bib}
\input{templates/cfg/zz_config_bib}

\usepackage{indentfirst}

\newcommand{\cmd}[1]{\texttt{#1}}

% Alternate row colours in every table (xcolor is loaded with the table option by the layout)
% Breathing room between tables and the running text
\setlength{\intextsep}{16pt plus 2pt minus 2pt}
\setlength{\textfloatsep}{18pt plus 2pt minus 4pt}
\setlength{\floatsep}{16pt plus 2pt minus 2pt}

\definecolor{rowtint}{gray}{0.93}
\AtBeginEnvironment{table}{\global\csname @rowcolorstrue\endcsname\rowcolors{2}{rowtint}{white}}
\newcolumntype{P}[1]{>{\raggedright\arraybackslash}p{#1}}
\setcounter{secnumdepth}{2}
\newcommand{\todo}[1]{\textcolor{red}{[TODO: #1]}}

\begin{document}

%=====================================================
% Title page and abstract
%=====================================================

\input{tex/\ttitlepage}

\begin{center}
\begin{minipage}{0.86\textwidth}
\section*{Abstract}

The earlier reports asked whether a local open-weight model can act as a coding worker. This fourth report asks
how it should be used: alone, or as the implementing layer (L2) under a cloud planner (L1) that writes work
packages, reviews the result and retries or takes over on failure. Two local models served by Strata on one
workstation (Gyro-L, Gyro-M) are compared with six cloud models on function-level sets, five small changes from a
real repository and two larger multi-file changes, alone and in L1/L2 pairs. Four findings stand out. First, on
small single-file changes the local Gyro-L alone solves 14 of 15 at no cloud cost; a planner on top adds at most a
few points and costs more than Sonnet alone in single-task groups. Second, on larger multi-file changes Sonnet as
L1 over Gyro-L as L2 solves 4 of 4, like Sonnet alone, at 35\,\% less cloud cost but about six times the wall time;
Gyro-L alone solves 3 of 4. Third, Sonnet over Haiku solves everything at 17\,\% less cloud cost and is faster than
Sonnet alone. Fourth, the planner's waiting is a cost of its own: an L1 re-reads 35--80 thousand cached tokens per
minute while its workers run, and one blocking wait per worker instead of periodic checks lowers its cost by about
11\,\% at equal pass rate. We recommend a balanced default (local worker, cloud takeover), the hierarchy for large
non-urgent work, Haiku as L2 when speed matters, and one switch between these regimes.

\end{minipage}
\end{center}

\clearpage

%=====================================================
% Table of contents and list of tables
%=====================================================

\sequals{\showpageheaders}{true}{
	\pagestyle{fancy}
	\fancyhead{}
	\fancyhead[RO,LE]{\thepage}
	\fancyhead[LO]{\slshape CONTENTS}
	\fancyhead[RE]{\slshape CONTENTS}
}{}
\sequals{\showpagefooters}{true}{\fancyfoot[C]{}}{}
\pdfbookmark[-1]{\contentsname}{toc}
\tableofcontents
\addtocontents{lot}{\protect\vspace{10pt}}
\listoftables
\clearpage

\sequals{\showpageheaders}{true}{
	\pagestyle{fancy}
	\fancyhead{}
	\fancyhead[LE,RO]{\thepage}
	\fancyhead[RE]{\slshape \rightmark}
	\fancyhead[LO]{\slshape \rightmark}
}{}
\sequals{\showpagefooters}{true}{\fancyfoot[C]{\pfooter}}{}
\renewcommand{\sectionmark}[1]{\markboth{\thesection.\ #1}{\thesection.\ #1}}
\renewcommand{\subsectionmark}[1]{}
\clearpage

%=====================================================
% Body
%=====================================================

\section{Introduction}\label{sec:intro}

The first three reports established that open-weight models on one workstation can finish real coding tasks once
the serving stack, the prompt shape and the harness are right, and identified which models and quantisations do it
best. They measured the local model as the only worker. In daily use, however, the agent setup is hierarchical: an
orchestrating session (L0) hands work to planners (L1), and planners hand concrete changes to implementing
workers (L2). The questions of this report are therefore economic and structural. When does a local model alone
suffice? When does a cloud planner on top pay for itself? Which planner and which worker, and at what cost in wall
time? All workers were started through PAI \citep{Nott-2026aa} with Claude Code as the agent harness
\citep{Anthropic-2026aa}.

\section{Setting}\label{sec:setting}

All local runs were served by Strata v0.2.13 with a Metal port on one Apple-silicon workstation (128\,GB unified
memory), one model at a time on a single port, without reasoning (reasoning effort none), with speculative
decoding through an MTP head (draft acceptance about 90\,\%). Measured pure decode speed at the server was
91~tokens/s for Gyro-L and 112~tokens/s for Gyro-M.

\begin{table}[htbp]
\centering
\small
\caption{Models of this report}\label{tab:models4}
\begin{tabular}{P{3.6cm}P{2.6cm}P{7.4cm}}
\toprule
Model & Placement & Role \\
\midrule
Gyro-L & local & Worker (L2) and stand-alone; pair-trellis experts, separate per-layer embedding table \\
Gyro-M & local & Worker (L2) and stand-alone \\
Claude Sonnet \citep{Anthropic-2026ab} & cloud & Planner (L1), stand-alone yardstick \\
Claude Opus & cloud & Stand-alone yardstick \\
Claude Haiku & cloud & Worker (L2) \\
GLM, GLM-fast & cloud & Planner and worker pair; stand-alone \\
Kimi & cloud & Planner and worker pair; stand-alone \\
\bottomrule
\end{tabular}
\end{table}

\section{Method}\label{sec:method}

\subsection{Task sets}

\emph{short}: 20 questions with exact answers about code. \emph{easy} and \emph{hard}: 40 function-level Python
tasks each, graded by hidden tests. \emph{real (small)}: five changes taken from the PAI repository (bug fixes and
small features touching one to three files), graded by hidden behaviour tests derived from the original commits and
restricted to what the task text states. \emph{real (large)}: two multi-file changes from the same repository
(session identifiers in notes across the session and router modules; handling of renamed notes in the storage
registry), each requiring two to four coordinated edits. For every task the reference solution passes its tests
and the unchanged starting code fails them; tasks do not depend on the state of the machine.

\subsection{Hierarchy}

Related tasks are grouped and each group gets one planner (L1): the small tasks form three groups (worker
configuration, build tooling, sessions), each large task its own group. Per task the L1 reads only what it needs,
writes a self-contained work package (cause, exact target, required behaviour with edge cases, the proof to bring)
and starts one L2 on its own clone. It reviews only the diff and the visible test result, allows one retry with a
sharper package, and takes over with the smallest correct edit if the retry fails. Two waiting styles were
compared: \emph{polling} (periodic checks of the worker) and \emph{blocking} (all workers dispatched first, then
one long blocking wait per worker).

\subsection{Measures}

Pass rate on the hidden tests; retries and takeovers; wall time per run; cloud cost as list price of the cloud
layers (the local layer costs nothing), and for providers with other price schedules (GLM, Kimi) tokens per layer.
A loop guard stops a worker that repeats an identical write or command; a stopped run keeps its result.

\section{Results}\label{sec:results}

\subsection{Models alone}

Table~\ref{tab:alone} gives the share of tasks each model solves on its own, Table~\ref{tab:speed} the time it
needs per task and its effective output speed.

\begin{table}[H]
\centering
\small
\caption{Models alone: share of tasks solved}\label{tab:alone}
\begin{tabular}{llrrrr}
\toprule
Model & Where & easy & hard & short & real \\
\midrule
Gyro-L & local & 95.0\,\% & 86.7\,\% & 92.0\,\% & 93\,\% \\
Gyro-M & local & 97.5\,\% & 76.7\,\% & 81.0\,\% & 80\,\% \\
Sonnet & cloud & 100\,\% & 100\,\% & 100\,\% & 100\,\% \\
Opus & cloud & 100\,\% & 100\,\% & 100\,\% & 100\,\% \\
GLM & cloud & 99.5\,\% & 98.0\,\% & 100\,\% & 84\,\% \\
GLM-fast & cloud & 100\,\% & 90.5\,\% & 99\,\% & 96\,\% \\
Kimi & cloud & 98.3\,\% & 95.0\,\% & 100\,\% & -- \\
\bottomrule
\end{tabular}
\par\smallskip{\footnotesize\raggedright easy, hard and short are the function-level and question sets; real is the five small repository changes.\par}
\end{table}

\begin{table}[H]
\centering
\small
\caption{Models alone: speed}\label{tab:speed}
\begin{tabular}{lrrrrrr}
\toprule
 & \multicolumn{4}{c}{Seconds per task} & \multicolumn{2}{c}{Tokens/s} \\
\cmidrule(lr){2-5}\cmidrule(lr){6-7}
Model & easy & hard & short & real & sets & real \\
\midrule
Gyro-L & 3.6 & 11.0 & 2.3 & 369 & 63 & 25 \\
Gyro-M & 5.5 & 21.2 & 1.6 & 207 & 80 & 43 \\
Sonnet & 4.5 & 5.7 & 3.7 & 130 & 100 & 39 \\
Opus & 5.1 & 6.2 & 4.7 & 171 & 80 & 48 \\
GLM & 9.3 & 15.3 & 7.2 & 346 & 53 & 51 \\
GLM-fast & 10.9 & 23.8 & 4.6 & 490 & 43 & 39 \\
Kimi & 7.0 & 13.1 & 5.4 & -- & 38 & -- \\
\bottomrule
\end{tabular}
\par\smallskip{\footnotesize\raggedright Seconds per task: mean wall time for one task of each set, from start to graded result. Tokens/s: output tokens per second of wall time over the three sets together and over the real changes; it includes waiting for tools and is therefore lower than the pure decode speed.\par}
\end{table}

Gyro-L is the stronger local worker on hard and real tasks; Gyro-M decodes faster but fails more of the harder
tasks. The cloud models other than GLM solve the small repository changes completely. On the real changes the
local models need 1.6 (Gyro-M) to 2.8 (Gyro-L) times the wall time of Sonnet per task.

\subsection{Hierarchy on small changes}

Table~\ref{tab:small} compares the planner and worker pairs on the five small changes with Sonnet and with
Gyro-L working alone.

\begin{table}[H]
\centering
\small
\caption{Hierarchy on five small repository changes}\label{tab:small}
\resizebox{\textwidth}{!}{\begin{tabular}{lrrlrr}
\toprule
Setup & Runs & Passed & Retries / takeovers & Cloud per run & Wall per run \\
\midrule
Sonnet alone & 5 & 25/25 & -- & \$0.72 & 11.5\,min \\
Gyro-L alone & 3 & 14/15 & -- & \$0 & 30.9\,min \\
Sonnet $\rightarrow$ Gyro-L, polling & 5 & 22/24 & 2 / 1 & \$0.61 ($-15$\,\%) & 22.3\,min \\
Sonnet $\rightarrow$ Gyro-L, blocking & 3 & 15/15 & 1 / 1 & \$0.56 ($-22$\,\%) & 27.1\,min \\
Sonnet $\rightarrow$ Gyro-M & 3 & 13/14 & 0 / 0 & \$0.53 ($-26$\,\%) & 21.5\,min \\
Sonnet $\rightarrow$ Haiku & 3 & 15/15 & 1 / 0 & \$0.60 ($-17$\,\%) & 8.5\,min \\
GLM $\rightarrow$ GLM-fast & 3 & 15/15 & 0 / 0 & 135k out, 9.8M cache read & 91\,min \\
Kimi $\rightarrow$ Kimi & 3 & 15/15 & 1 / 0 & 32k out, 4.4M cache read & 40\,min \\
\bottomrule
\end{tabular}}
\par\smallskip{\footnotesize\raggedright Three groups of related changes, one planner each. Cloud cost per run compared with Sonnet alone. GLM and Kimi are given in planner tokens per run; a Sonnet planner uses about 16 thousand output and 0.9 million cache-read tokens per run.\par}
\end{table}

Every pair with Sonnet as planner costs less than Sonnet alone. The saving comes from the two groups that hold
two related changes each; in the group with a single change the pairs cost more than Sonnet alone on that change
(\$0.10--0.11 against \$0.07), because the planner's understanding of the code is not shared with a second
task. Against Gyro-L working alone the planner adds at most one solved change in fifteen.

\subsection{Hierarchy on larger changes}

Table~\ref{tab:large} shows the two larger multi-file changes, each handled by its own planner.

\begin{table}[H]
\centering
\small
\caption{Hierarchy on two larger multi-file changes}\label{tab:large}
\begin{tabular}{lrrlrr}
\toprule
Setup & Runs & Passed & Retries / takeovers & Cloud per run & Wall per run \\
\midrule
Sonnet alone & 2 & 4/4 & -- & \$0.77 & 4.3\,min \\
Gyro-L alone & 2 & 3/4 & -- & \$0 & 29.9\,min \\
Sonnet $\rightarrow$ Gyro-L, blocking & 2 & 4/4 & 2 / 0 & \$0.50 ($-35$\,\%) & 25.4\,min \\
\bottomrule
\end{tabular}
\end{table}

Here the planner pays: both retries turned a failed first attempt of the local worker into a pass, the pair
matches Sonnet alone at a third less cloud cost, and it beats the local worker alone. The price is wall time, about
six times that of Sonnet alone.

\subsection{Cost of the planner}

While its workers run, a planner re-reads its cached context on every turn: 35 to 80 thousand cached tokens per
minute of wall time, for every planner model measured. A slow worker therefore makes the planner expensive. With
blocking waits the planner's cost fell by about 11\,\% at equal pass rate (Table~\ref{tab:small}, polling against
blocking), and its cached reads per minute fell from 47--59 to 24--40 thousand. GLM and Kimi as planners reach full
quality but need two to eight times Sonnet's tokens and two to four times its wall time. In this setup GLM no longer
earns a place: alone it solves fewer real changes than every model but Gyro-M (84\,\%), and as a planner it is the
slowest and most token-hungry option (91\,min and about 9.8 million cache-read tokens per run, against 22\,min and
0.9 million for Sonnet over Gyro-L).

\section{Discussion}\label{sec:discussion}

The value of a planner depends on the size of the work. On small, well-described changes the local worker alone is
the efficient choice; a planner spends about as much understanding the change as a cloud model spends doing it.
On larger work, where a single model loses track, the planner's diagnosis and its one targeted retry restore full
quality while the implementation tokens stay local. The trade-off is time: the local worker is several times slower
than a cloud model doing the work itself. Haiku as worker inverts the trade-off: full quality, faster than Sonnet
alone and still cheaper, at the cost of doing no work locally. Grouping related work under one planner and making
the planner wait without polling are the two structural levers; both are cheap to enforce in the orchestration
layer.

\section{Recommendations}\label{sec:recommendations}

The results suggest three layers with distinct jobs (Table~\ref{tab:roles}) and a route that depends on the size
of the work and on whether time or cloud budget matters more (Table~\ref{tab:modes}).

\begin{table}[H]
\centering
\small
\caption{Layers and their jobs}\label{tab:roles}
\begin{tabular}{P{1.2cm}P{3.2cm}P{9.2cm}}
\toprule
Layer & Model & Job \\
\midrule
L0 & Opus (the operator's session) & Understands the request, groups related work, chooses the route by size \\
L1 & Sonnet & Plans a group of related changes: work packages, review of diff and test result, one retry, takeover \\
L2 & Gyro-L (local), or Haiku & Implements one work package \\
\bottomrule
\end{tabular}
\end{table}

\begin{table}[H]
\centering
\small
\caption{Route by mode and size of the work}\label{tab:modes}
\begin{tabular}{P{1.7cm}P{4.6cm}P{5.0cm}P{2.4cm}}
\toprule
Mode & Single small change & Several related or one larger multi-file change & On failure \\
\midrule
speed & L0 $\rightarrow$ Sonnet & L0 $\rightarrow$ Sonnet (L1) $\rightarrow$ Haiku (L2) & Sonnet takes over \\
balanced & L0 $\rightarrow$ Gyro-L (L2), no planner & L0 $\rightarrow$ Sonnet (L1) $\rightarrow$ Gyro-L (L2) & Sonnet takes over \\
budget & L0 $\rightarrow$ Gyro-L (L2), no planner & L0 $\rightarrow$ Sonnet (L1) $\rightarrow$ Gyro-L (L2) & Sonnet only after a failed retry \\
\bottomrule
\end{tabular}
\par\smallskip{\footnotesize\raggedright speed favours wall time (Sonnet alone 11.5\,min, Sonnet over Haiku
8.5\,min per run of five small changes); balanced and budget keep the implementation local (Sonnet over Gyro-L
\$0.50 instead of \$0.77 on the larger changes). balanced is the default.\par}
\end{table}

Beyond the route, the planner should follow three rules: group related work under one planner, wait for its
workers with one blocking call each, and review only the diff and the test result. A planner is not started for a
single small change.

These recommendations are implemented in PAI \citep{Nott-2026aa}. The regime is one setting,
\cmd{pai worker mode speed|balanced|budget}, persisted in the worker configuration and shown in the status line;
\emph{balanced} is the default. A session marked urgent (\cmd{/now}) runs as \emph{speed} while the mark is
set. Orchestration requests for the same repository that arrive within a short window (20\,s by default) are
merged into one planner job; planners wait for their workers with one blocking call each and review only the diff
and the test result; a loop guard stops a local worker after three identical writes.

\section{Limitations}\label{sec:limitations}

The larger-change result rests on two tasks with two runs each; five to eight further multi-file changes are
needed before a size threshold for delegation is fixed. All repository tasks come from one TypeScript code base.
Local and cloud runs partly overlapped in time and shared the harness's test runner, which lengthens wall times
slightly for both. Cloud costs are list prices; under a subscription they translate into quota. Local models ran
without reasoning; a reasoning configuration was not measured.

\clearpage
\markboth{\MakeUppercase{\slshape References}}{}
\phantomsection
\addcontentsline{toc}{section}{References}
\printbibliography[title={References}]

\end{document}
